\documentclass[conference]{IEEEtran}
\IEEEoverridecommandlockouts
\usepackage{cite}
\usepackage{amsmath,amssymb,amsfonts}
\usepackage{graphicx}
\usepackage{textcomp}
\usepackage{xcolor}
\usepackage{booktabs}
\usepackage{siunitx}
\usepackage{multirow}
\usepackage{enumitem}
\usepackage{url}
\usepackage[colorlinks=false, hidelinks]{hyperref}
\usepackage{pgfplots}
\pgfplotsset{compat=1.18}

\def\BibTeX{{\rm B\kern-.05em{\sc i\kern-.025em b}\kern-.08em
    T\kern-.1667em\lower.7ex\hbox{E}\kern-.125emX}}

\begin{document}

\title{Graph Context, Not More Context: Isolating What\\Repository Structure Adds to LLM Code Review}

\author{
\IEEEauthorblockN{Abdelrahman (Abdu) Khattab}
\IEEEauthorblockA{\textit{Hasso Plattner Institute, University of Potsdam}\\
Potsdam, Germany \\
abdelrahman.khattab@student.hpi.de}
% TODO: confirm co-author list and affiliations before submission.
}

\maketitle

\begin{abstract} 
Large language models review pull requests from the diff alone, which makes
them structurally blind to defects whose consequences lie outside the changed
file. A natural remedy is to supply repository context, but the literature
tends to treat ``more context'' as a single lever and reports aggregate
quality gains that cannot distinguish a real capability from a general
verbosity effect. We separate the two. We compare a diff-only baseline
against three context conditions --- a repository knowledge graph, embedding
retrieval, and their combination --- under two complementary designs: an
observational study of 40 merged pull requests from five open-source
repositories, scored by a three-judge LLM panel on a 25-criterion rubric,
and a controlled injection study of 40 synthetic defects with an objective
static-structural oracle. The observational study shows modest aggregate
gains that are easy to over-read; its diagnostic value lies in \emph{where}
the gains land. Graph context concentrates its wins in the nine
dependency-related criteria (net $+24$) and is flat elsewhere (net $+1$),
whereas retrieval spreads its wins evenly across both bands ($+17$ and
$+18$). The injection study makes the consequence unambiguous: on 28 defects
whose effect is provably outside the changed file, the diff-only baseline
detects $0$, embedding retrieval detects $1$, and graph context detects $15$
($\Delta = +0.54$, 95\% CI $[+0.36, +0.71]$, $p<0.0001$). A placebo band of
12 defects visible inside the diff shows no graph advantage
($\Delta = -0.08$), which rules out a context-volume explanation. An
idealised graph reaches $26/28$, locating most of the remaining gap in graph
construction rather than in the language model. We conclude that the useful
variable is the \emph{shape} of the retrieved context, not its quantity, and
that cross-file defect detection is the capability worth measuring.
\end{abstract}

\begin{IEEEkeywords}
code review, large language models, knowledge graphs,
retrieval-augmented generation, program analysis, empirical evaluation
\end{IEEEkeywords}

\section{Introduction} 
\label{sec:intro}

Modern code review is expensive, and it is expensive in a specific way:
reviewers spend much of their effort establishing what a change will affect
beyond the lines they are reading~\cite{bacchelli2013, sadowski2018}. The
comments practitioners rate as most useful are disproportionately the ones
that connect a change to consequences elsewhere in the
system~\cite{bosu2015}. Large language models have made automated review
comments fluent enough to be worth deploying~\cite{li2022codereviewer,
tufano2022codereview, fan2023llmse}, but the standard deployment feeds the
model a diff and nothing else. Whatever the model knows about the rest of the
repository it must have memorised during pre-training.

This produces a blind spot that is structural rather than incidental. If a
function signature changes and its callers live in other files, a diff-only
reviewer has no mechanism for noticing them. The obvious response is to
retrieve additional repository context, and two families of technique are
available: embedding-based retrieval, which finds text that resembles the
diff~\cite{lewis2020rag, zhang2023repocoder}, and graph-based
representations, which encode the relations a compiler would
resolve~\cite{allamanis2018programgraphs, guo2020graphcodebert,
liu2024graphcoder, yang2025kgcompass}.

Our concern is with how the benefit of such context is usually established.
The common protocol reports an aggregate quality score before and after
adding context. Such a score cannot separate three very different
explanations for an improvement: the model acquired a capability it
previously lacked; the model produced longer output and thereby satisfied
more rubric items; or the judge rewarded the appearance of technical
specificity. The distinction matters practically, because only the first
explanation predicts that the system will catch a defect a human reviewer
would have missed.

We therefore pair an observational study with a controlled one. The
observational study (Experiment 1) follows the conventional protocol closely
enough to be comparable to prior work: 40 merged pull requests, four context
conditions, a three-judge LLM panel, a 25-criterion rubric derived from the
code review literature. We report its aggregate effects, but we treat them as
weak evidence and use the study primarily as a diagnostic, by pre-declaring
which nine criteria the graph context is supposed to affect and treating the
remaining sixteen as a negative control band. The controlled study
(Experiment 2) removes the judge from the ground truth entirely: we inject
defects whose consequences are resolved statically, so that ``did the
reviewer find it'' has an answer that does not depend on an LLM's opinion.

The two studies disagree in a way we think is the paper's main finding. In
aggregate the context conditions look interchangeable and unexciting --- the
graph and retrieval conditions differ by a few tenths of a point on a
25-point scale. On cross-file defect detection they are not remotely
interchangeable: $15/28$ against $1/28$. An evaluation that stops at the
aggregate score would have concluded that the choice between them barely
matters.

Our contributions are:

\begin{itemize}[leftmargin=*]
\item A within-subject comparison of graph and embedding context for LLM code
      review on 40 real merged pull requests, reported with paired confidence
      intervals and a pre-declared criterion split that separates targeted
      from generic gains (Section~\ref{sec:exp1}).
\item A controlled defect-injection methodology for code review context that
      replaces judge opinion with a static-structural oracle, together with a
      placebo band that distinguishes a structural capability from a
      context-volume artefact (Section~\ref{sec:exp2}).
\item Evidence that the two context types act through different mechanisms,
      and that the graph advantage is bounded mainly by graph construction
      rather than by the language model: an idealised resolver reaches
      $26/28$ where the deployed one reaches $15/28$
      (Section~\ref{sec:discussion}).
\end{itemize}

\section{Background and Related Work}
\label{sec:related}

\subsection{What makes a review comment useful}
Studies of review practice consistently find that reviewers' central
difficulty is understanding a change's context rather than reading its
syntax~\cite{bacchelli2013, sadowski2018}. Bosu et al.\ characterise useful
comments at Microsoft and find that comments identifying consequences beyond
the immediate change carry disproportionate
value~\cite{bosu2015}. Recent work continues to refine what reviewers
actually act on, including the role of explanations~\cite{widyasari2025explanations}
and of concrete suggestions~\cite{bouraffa2025suggestions}. We use this
literature in two ways: to justify the rubric, and to justify our choice of
cross-file consequence detection as the capability worth isolating.

\subsection{Automated review comment generation and its evaluation}
Pre-trained and instruction-tuned models have been applied to review comment
generation with steadily improving fluency~\cite{li2022codereviewer,
tufano2022codereview, fan2023llmse}. Evaluation has not kept pace. Hong
et al.\ show that common benchmarks and reference-similarity metrics
substantially misrank code review comment generators, and argue for
evaluation that targets what reviewers need~\cite{hong2024deepcrceval}. LLM
panels are now a standard substitute for human judgement, with known
caveats around position and verbosity bias~\cite{zheng2023judging}. Our
Experiment 1 inherits exactly these weaknesses, which is why we do not rest
the paper's claim on it.

\subsection{Repository context: text similarity versus program structure}
Embedding retrieval extends a model's effective context by surfacing
similar text~\cite{lewis2020rag}, and works well for repository-level code
completion when the needed code resembles the code being
written~\cite{zhang2023repocoder}. Structural representations take the
opposite route, encoding relations that a compiler resolves --- calls,
imports, inheritance, data flow~\cite{aho2006compilers,
allamanis2018programgraphs, guo2020graphcodebert}. GraphCoder shows that
code-context graphs improve completion~\cite{liu2024graphcoder}, and
KGCompass uses a knowledge graph for issue-localisation
retrieval~\cite{yang2025kgcompass}. The comparison we make is between these
two families as \emph{context suppliers to a reviewer}, holding the reviewing
model fixed. To our knowledge the specific question --- whether the two are
substitutes or complements for cross-file defect detection --- has not been
answered with an objective oracle.

\subsection{Relation to the pull request assessment pipeline}
Our pull request dataset and the surrounding tooling derive from Mariotto
et al.'s work on assessing and enhancing pull requests at
scale~\cite{mariotto2025, mariotto2025esem}, which aligns review output with
developer competencies; related work in the same line studies the orders in
which reviews are usefully organised~\cite{bouraffa2025orders}. That line of
work asks how to shape review output for a reader. We ask a prior question:
what does the reviewer need to be able to see at all.
% TODO: citations still needed for (a) the Joern code property graph and the
% original CPG formulation, (b) GraphRAG-style graph retrieval, (c) mutation
% testing as the ancestor of our injection design. Do not invent keys.

\section{Approach}
\label{sec:approach}

We hold the reviewing model, the prompt scaffold, and the scoring procedure
fixed, and vary only the context block that is inserted into the prompt. Four
conditions are compared throughout.

\textbf{baseline} receives the pull request title, description, and diff, and
nothing else. This is the prevailing deployment.

\textbf{kg} additionally receives a rendered neighbourhood of the repository
at the pull request's head commit, naming the changed files, the files
reported as depending on them, the relation type by which each is reached, and
the associated tests. The deployed builder for this condition is deliberately
cheap, and we state its construction plainly because the term \emph{knowledge
graph} invites a stronger reading: dependents are found by searching the
repository for the basename of each changed file, filtering hits by source
extension, and capping the result at fifteen per file. Nothing is parsed and no
relation is verified, so the block over-collects and can label a file as a
dependent on the strength of a mention in a comment or docstring. Any effect it
produces is an effect of this construction rather than of graph-structured
context in general, which is why we re-run the condition with a builder that
resolves relations properly (Section~\ref{sec:exp1-builder}) and why the causal
experiment uses that builder instead.

\textbf{rag} additionally receives the eight highest-scoring chunks retrieved
from the repository by embedding similarity. Files near the changed paths are
chunked at 2\,000 characters with a 200-character overlap and embedded with
\texttt{text-embedding-3-small}; the queries are the leading 4\,000 characters
of each changed file rather than the diff itself, and chunks belonging to the
changed files are excluded, so retrieval contributes only material from outside
the diff. This is the conventional retrieval-augmented configuration.

\textbf{hybrid} receives both blocks.

Two properties of this design are worth stating because they carry the
paper's internal validity. First, the conditions are nested rather than
disjoint: every condition contains the baseline's information, so a
condition can only lose by \emph{diluting} what the model attends to, never
by withholding the diff. Second, the context blocks are of broadly
comparable size, which is what allows the injection study's placebo band to
function as a volume control (Section~\ref{sec:exp2}).

\subsection{Graph construction}
Three builders appear in this paper, and the differences between them are
themselves a result. The first is the lexical one described above: it matches
filenames, is language-agnostic and cheap, and buys recall at the cost of
precision. The second parses each file's abstract syntax tree and matches
resolved module paths, which removes the spurious relations but misses those
requiring type resolution. The third uses a code property graph, resolving
calls, data flow and inheritance across files at the cost of a
language-specific frontend and a much higher build time. Experiment 1's
headline uses the lexical builder, because that is the configuration the study
was run and pre-registered with; we report the code property graph as a
builder-sensitivity analysis within Experiment 1
(Section~\ref{sec:exp1-builder}) and use it as the deployed builder in
Experiment 2. The two experiments therefore do not share a graph
construction, and we are careful not to let the causal result of the second
vouch for the builder used in the first.
% TODO: cite Joern / CPG formulation here once the key is added to the bib.

\section{Experiment 1: Observational Study}
\label{sec:exp1}

\subsection{Design}
We sample 40 merged pull requests from five open-source repositories
(Grafana 13, Apache Kafka 8, scikit-learn 8, Godot 6, Jenkins 5), spanning
Go, Java, Python, C++, and TypeScript. For each pull request we generate one
review under each of the four conditions, giving 160 reviews. Each review is
scored against 25 binary criteria by three LLM judges (gpt-4o-mini, gpt-4o,
gemini-2.5-flash), and the majority of three is the cell of record, with ties
resolved to zero. This yields $160 \times 25 \times 3 = 12\,000$ judged cells
and $4\,000$ cells of record. Table~\ref{tab:config} gives the full
configuration for both experiments.

% Generated by scripts/generate_paper_tables.py — do not edit by hand.
% Source: results/checklist_evaluation_llm_multi__v2.json + injection manifest
\begin{table}[t]
\centering
\caption{Configuration of both experiments. Generation temperatures differ
because the two pipelines call different entry points; within an experiment
every arm shares one setting, so all contrasts are internally valid. No
generation seed is sent to the model, so Experiment 1's determinism rests on
$T=0$.}
\label{tab:config}
\footnotesize
\setlength{\tabcolsep}{3.5pt}
\begin{tabular}{lcc}
\toprule
& Experiment 1 & Experiment 2 \\
\midrule
Unit & merged PR & injected defect \\
$n$ & 40 & 40 (28 struct., 12 control) \\
Repositories & 5 & 3 \\
Generator & gpt-4o & gpt-4o \\
Temperature & 0.0 & 0.3 \\
Diff cap (chars) & 50\,000 & 12\,000 \\
Judges & \multicolumn{2}{c}{gpt-4o-mini, gpt-4o, gemini-2.5-flash} \\
Aggregation & \multicolumn{2}{c}{majority of three, ties to 0} \\
Outcome & rubric /25 & defect detected \\
Ground truth & LLM panel & static oracle \\
Statistics seed & \multicolumn{2}{c}{2026} \\
\bottomrule
\end{tabular}
\end{table}


The rubric's 25 criteria are grounded in the review literature
\cite{bacchelli2013, bosu2015, sadowski2018} and in ISO/IEC 25010.
% TODO: add the ISO/IEC 25010 bib entry.
Before looking at any scores we designated nine of them --- those concerning
cross-file impact, dependency identification, and downstream consequence ---
as the \emph{dependency subset}, and the remaining sixteen as a negative
control band. This split is the study's main diagnostic instrument: graph
context has a mechanism for improving the first group and no particular
mechanism for improving the second.

Uncertainty is reported by percentile bootstrap over pull requests
($B=10\,000$) and paired sign-flip permutation tests
($B=20\,000$), both with seed 2026. All comparisons are within-subject: each
condition is compared against the baseline review of the \emph{same} pull
request.

\subsection{Aggregate results}
Table~\ref{tab:exp1-means} reports per-condition means and
Table~\ref{tab:exp1-deltas} the paired differences.

% Generated by scripts/generate_paper_tables.py — do not edit by hand.
% Source: results/BOOTSTRAP_STATS_v2.json
\begin{table}[t]
\centering
\caption{Experiment 1: mean rubric score per condition with 95\% percentile
bootstrap confidence intervals ($B=10\,000$, seed 2026), $n=40$ pull
requests. \emph{baseline} is diff-only and \emph{hybrid} is kg+rag. Total is
out of 25 criteria; the dependency subset is the nine pre-registered criteria
the graph context is designed to affect.}
\label{tab:exp1-means}
\footnotesize
\setlength{\tabcolsep}{3.5pt}
\begin{tabular}{lcc}
\toprule
Condition & Total /25 & Dependency /9 \\
\midrule
baseline & 9.20 [8.70, 9.70] & 4.97 [4.67, 5.28] \\
kg & 9.82 [9.28, 10.43] & 5.58 [5.22, 5.95] \\
rag & 10.07 [9.57, 10.60] & 5.40 [5.00, 5.78] \\
hybrid & 9.93 [9.47, 10.38] & 5.50 [5.22, 5.78] \\
\bottomrule
\end{tabular}
\end{table}

% Generated by scripts/generate_paper_tables.py — do not edit by hand.
% Source: results/BOOTSTRAP_STATS_v2.json
\begin{table*}[t]
\centering
\caption{Experiment 1: paired differences against the diff-only baseline on
the same pull request. Two-sided paired sign-flip permutation test
($B=20\,000$, seed 2026); $d_z$ is the paired Cohen's $d$.
\textsuperscript{*}$p<.05$, \textsuperscript{**}$p<.01$,
\textsuperscript{***}$p<.001$. Each context type reaches significance on the
metric it targets, and the hybrid on both.}
\label{tab:exp1-deltas}
\begin{tabular}{lccccccc}
\toprule
& \multicolumn{3}{c}{Total /25} & \multicolumn{3}{c}{Dependency /9} \\
\cmidrule(lr){2-4} \cmidrule(lr){5-7}
Mode vs baseline & $\Delta$ [95\% CI] & $p$ & $d_z$ & $\Delta$ [95\% CI] & $p$ & $d_z$ \\
\midrule
kg & +0.62 [+0.05, +1.20] & 0.0537 & +0.33 & +0.60 [+0.23, +0.97] & 0.0066\textsuperscript{**} & +0.47 \\
rag & +0.88 [+0.30, +1.45] & 0.0073\textsuperscript{**} & +0.47 & +0.42 [+0.05, +0.82] & 0.0573 & +0.33 \\
hybrid (kg+rag) & +0.72 [+0.10, +1.35] & 0.0374\textsuperscript{*} & +0.36 & +0.53 [+0.17, +0.90] & 0.0083\textsuperscript{**} & +0.46 \\
\bottomrule
\end{tabular}
\end{table*}


Every context condition improves on the baseline, and every improvement is
small: between $+0.63$ and $+0.88$ points on a 25-point scale. Read as a
horse race, the result is uninformative --- retrieval wins on total score,
the graph wins on the dependency subset, the hybrid is in between, and the
confidence intervals overlap heavily. Notably, the graph condition's
total-score effect does not reach significance ($p = 0.054$) while its
dependency-subset effect does ($p = 0.0066$), and the pattern reverses for
retrieval. This is the first hint that the two conditions are not doing the
same thing, and it is the reason we do not report a single headline
percentage for ``adding context.''

\subsection{Where the gains land}
\label{sec:exp1-attribution}

Table~\ref{tab:exp1-attribution} decomposes each condition's gains into wins
and losses per criterion band. A win is a criterion the treatment satisfies
and the baseline does not, on the same pull request.

% Generated by scripts/generate_paper_tables.py — do not edit by hand.
% Source: results/checklist_evaluation_llm_multi__v2.json
\begin{table}[t]
\centering
\caption{Where the gains land (Experiment 1, $n=40$). A win is a criterion the
treatment satisfies and the baseline does not, on the same pull request. The
sixteen non-dependency criteria act as a built-in negative control: a generic
quality lift would spread evenly across both bands.}
\label{tab:exp1-attribution}
\footnotesize
\setlength{\tabcolsep}{3.5pt}
\begin{tabular}{lcccccc}
\toprule
& \multicolumn{3}{c}{Dependency (9 criteria)} & \multicolumn{3}{c}{Other (16 criteria)} \\
\cmidrule(lr){2-4} \cmidrule(lr){5-7}
Mode & W & L & net & W & L & net \\
\midrule
kg & 42 & 18 & \textbf{+24} & 34 & 33 & \textbf{+1} \\
rag & 38 & 21 & \textbf{+17} & 39 & 21 & \textbf{+18} \\
\bottomrule
\end{tabular}
\end{table}


The two conditions have almost the same total effect and clearly different
signatures. Graph context nets $+24$ on the nine dependency criteria and
$+1$ on the sixteen others: it changed what the reviewer could see about
structure and left everything else alone. Retrieval nets $+17$ and $+18$:
it improved the review broadly and roughly uniformly, which is what a
general fluency or specificity lift looks like. The negative control band
behaves as a control band should for one condition and not for the other.

We stress how invisible this is at the aggregate level. The two conditions
differ by $0.25$ points in total score, well inside the confidence interval;
they differ by a factor of twenty-four in how selectively they act.

\subsection{Builder sensitivity}
\label{sec:exp1-builder}

Because the graph condition's effect depends on how many true relations the
builder resolves, we re-ran the graph and baseline arms with the code
property graph builder on the same rubric and the same judge panel. The code
property graph frontend cannot parse the five Go pull requests, so this
comparison uses the 35 remaining ones. Table~\ref{tab:exp1-builder} reports
both.

% Generated by scripts/generate_paper_tables.py — do not edit by hand.
% Source: results/BOOTSTRAP_STATS_v2.json + results/BOOTSTRAP_STATS_joern.json
\begin{table*}[t]
\centering
\caption{Builder sensitivity within Experiment 1: the graph arm against the
same diff-only baseline, under two graph construction methods, scored on the
\emph{same} 25-criterion rubric by the \emph{same} three-judge panel. The
code-property-graph run excludes five Go pull requests its frontend cannot
parse, hence $n=35$. A stronger resolver leaves the dependency effect broadly
similar but moves the total-score effect from inconclusive to significant.}
\label{tab:exp1-builder}
\begin{tabular}{llcccccc}
\toprule
& & \multicolumn{3}{c}{Dependency /9} & \multicolumn{3}{c}{Total /25} \\
\cmidrule(lr){3-5} \cmidrule(lr){6-8}
Graph construction & $n$ & $\Delta$ [95\% CI] & $p$ & $d_z$ & $\Delta$ [95\% CI] & $p$ & $d_z$ \\
\midrule
Import/AST resolver & 40 & +0.60 [+0.23, +0.97] & 0.0066\textsuperscript{**} & +0.47 & +0.62 [+0.05, +1.20] & 0.0537 & +0.33 \\
Code property graph & 35 & +0.69 [+0.31, +1.09] & 0.0032\textsuperscript{**} & +0.58 & +1.17 [+0.54, +1.77] & 0.0016\textsuperscript{**} & +0.61 \\
\bottomrule
\end{tabular}
\end{table*}


The stronger builder roughly doubles the total-score effect ($+0.63$ to
$+1.17$) and moves it from inconclusive to significant, while the
dependency-subset effect changes comparatively little ($+0.60$ to $+0.69$).
The direction is what the mechanism predicts: resolving more true relations
gives the reviewer more to say, and it is the criteria outside the dependency
subset that had the most room to move. We report this as sensitivity rather
than as a headline, for two reasons: the sample differs ($n=35$), and only
the graph arm was rebuilt, so this file supports the baseline-versus-graph
contrast and no other.

\subsection{Reliability}
Pairwise inter-judge agreement on the 25-criterion rubric is
$\kappa = 0.72$ (gpt-4o-mini vs gpt-4o), $0.71$ (gpt-4o vs
gemini-2.5-flash), and $0.59$ (gpt-4o-mini vs gemini-2.5-flash): substantial
for two pairs and moderate for the third.\footnote{The Gemini judge returned
usable labels for $3\,836$ of $4\,000$ cells; the majority-of-three rule with
ties resolved to zero absorbs the remainder conservatively.}

This is respectable reliability, and we want to be careful about what it does
and does not buy. Agreement establishes that the judges are measuring
\emph{something} consistently; it does not establish that they are measuring
review usefulness. Three models drawn from an overlapping training
distribution can agree with each other and be wrong in the same direction,
and the biases documented for LLM panels --- verbosity and apparent
specificity in particular~\cite{zheng2023judging} --- are exactly the ones
that would inflate a context condition's score without any gain in
capability. High $\kappa$ cannot detect that failure mode; only ground truth
outside the panel can. Hence Experiment 2.

\section{Experiment 2: Controlled Defect Injection}
\label{sec:exp2}

\subsection{Design}
Experiment 1 leaves two gaps. Its ground truth is an LLM panel, and its
outcome is a rubric score rather than a consequence anyone cares about
directly. Experiment 2 closes both. The design was pre-registered before the
run, including the arms, the bands, the decision thresholds, and an explicit
rule against back-filling missing cases.

We inject 40 synthetic defects into three repositories --- scikit-learn
(Python), kafka-streams (Java), and grafana-data (TypeScript) --- in two
kinds of band. \emph{Structural} bands ($n=28$) place a defect whose
consequence is provably outside the changed file: the dependents are
resolved independently of the reviewer and recorded with evidence lines in
the manifest before any review is generated. \emph{Local control} bands
($n=12$) place a defect whose consequence is visible inside the diff itself.

The local control band is the design's load-bearing element. Every context
condition adds tokens to the prompt, so any observed gain might simply
reflect a longer, more thorough prompt. On local defects, extra repository
context is irrelevant by construction, so a condition that gains there is
gaining from volume rather than structure. We pre-registered a parity window
of $[-0.15, +0.15]$ for this band.

Defect sites were selected direction-blind, by structural fan-out, without
consulting any review. Detection is judged by the same three-model panel with
the same majority rule as Experiment 1, and separately validated against a
deterministic oracle that checks whether the review actually names the
affected dependent. Two arms are added beyond the four conditions: an
\emph{idealised} graph, which is given the statically resolved dependents
directly and therefore measures the ceiling available to a perfect builder,
and an \emph{inheritance-augmented} graph, a pre-declared repair for a
specific blind spot.

\subsection{Structural results}
% Generated by scripts/generate_paper_figures.py — do not edit by hand.
% Source: experiments/2026-07-05_injection_exp2/out/results.json
% Requires \usepackage{pgfplots} and \pgfplotsset{compat=1.18}.
\begin{figure*}[t]
\centering
\pgfplotsset{
  detbar/.style={
    xbar, width=0.47\textwidth, height=4.1cm,
    xmin=0, xmax=1.19, xtick={0,0.25,0.5,0.75,1.0},
    bar width=7pt, y=0.52cm,
    axis lines*=left, axis line style={gray!50},
    xmajorgrids, grid style={gray!22, line width=0.4pt},
    tick label style={font=\scriptsize},
    xlabel style={font=\scriptsize, yshift=2pt},
    title style={font=\small, align=center, yshift=1pt},
    error bars/x dir=both, error bars/x explicit,
    error bars/error bar style={gray!70!black, line width=0.6pt},
    error bars/error mark options={rotate=90, mark size=1.4pt, gray!70!black},
    ytick={0,1,2,3,4,5},
    enlarge y limits={abs=0.5},
    clip=false,
  },
}
\begin{tikzpicture}
\begin{axis}[detbar,
    yticklabels={baseline,rag,hybrid,kg,kg+inherit.,kg idealised},
    ytick style={draw=none},
    xlabel={detection rate},
    title={Structural defects ($n{=}28$)\\\scriptsize consequence \emph{outside} the changed file},
  ]
  \addplot+[draw=none, fill=blue!58!black!78] plot coordinates {
        (0.0000,0) +- (0.0000,0) -= (0.0000,0)
        (0.0357,1) +- (0.0714,0) -= (0.0357,0)
        (0.4286,2) +- (0.1786,0) -= (0.1786,0)
        (0.5357,3) +- (0.1786,0) -= (0.1786,0)
        (0.7500,4) +- (0.1429,0) -= (0.1786,0)
        (0.9286,5) +- (0.0714,0) -= (0.1071,0)
  };
      \node[anchor=west,font=\scriptsize,gray!60!black] at (axis cs:0.0300,0) {0/28};
      \node[anchor=west,font=\scriptsize,gray!60!black] at (axis cs:0.1371,1) {1/28};
      \node[anchor=west,font=\scriptsize,gray!60!black] at (axis cs:0.6371,2) {12/28};
      \node[anchor=west,font=\scriptsize,gray!60!black] at (axis cs:0.7443,3) {15/28};
      \node[anchor=west,font=\scriptsize,gray!60!black] at (axis cs:0.9229,4) {21/28};
      \node[anchor=west,font=\scriptsize,gray!60!black] at (axis cs:1.0300,5) {26/28};
\end{axis}
\begin{axis}[detbar, at={(0.5\textwidth,0)}, anchor=south west,
    yticklabels={,,},
    ytick style={draw=none},
    xlabel={detection rate},
    title={Local controls ($n{=}12$)\\\scriptsize consequence \emph{inside} the diff},
  ]
  \addplot+[draw=none, fill=gray!45] plot coordinates {
        (0.9167,0) +- (0.0833,0) -= (0.1667,0)
        (1.0000,1) +- (0.0000,0) -= (0.0000,0)
        (0.9167,2) +- (0.0833,0) -= (0.1667,0)
        (0.8333,3) +- (0.1667,0) -= (0.2500,0)
        (0.8333,4) +- (0.1667,0) -= (0.2500,0)
        (0.8333,5) +- (0.1667,0) -= (0.2500,0)
  };
      \node[anchor=west,font=\scriptsize,gray!60!black] at (axis cs:1.0300,0) {11/12};
      \node[anchor=west,font=\scriptsize,gray!60!black] at (axis cs:1.0300,1) {12/12};
      \node[anchor=west,font=\scriptsize,gray!60!black] at (axis cs:1.0300,2) {11/12};
      \node[anchor=west,font=\scriptsize,gray!60!black] at (axis cs:1.0300,3) {10/12};
      \node[anchor=west,font=\scriptsize,gray!60!black] at (axis cs:1.0300,4) {10/12};
      \node[anchor=west,font=\scriptsize,gray!60!black] at (axis cs:1.0300,5) {10/12};
\end{axis}
\end{tikzpicture}
\caption{Defect detection by arm, with 95\% bootstrap confidence intervals and
raw counts. Left: defects whose consequence provably lies outside the changed
file. The diff-only baseline detects none, embedding retrieval detects one,
and graph context detects fifteen; the idealised resolver shows how much of
the remaining gap belongs to graph construction rather than to the model.
Right: the placebo band, where the defect is visible in the diff and no arm
should gain from added context --- and none does. The contrast between the two
panels is what distinguishes a structural capability from a context-volume
artefact.}
\label{fig:detection-ladder}
\end{figure*}

% Generated by scripts/generate_paper_tables.py — do not edit by hand.
% Source: experiments/2026-07-05_injection_exp2/out/results.json
\begin{table*}[t]
\centering
\caption{Experiment 2, structural bands ($n=28$): detection of an injected
defect whose consequence lies outside the changed file. Rates with 95\%
bootstrap confidence intervals; contrasts by paired permutation test
($B=20\,000$, seed 2026). The diff-only baseline never detects one.}
\label{tab:exp2-structural}
\begin{tabular}{lcccc}
\toprule
Arm & Detected & Rate [95\% CI] & $\Delta$ vs baseline [95\% CI] & $p$ \\
\midrule
baseline (diff only) & 0/28 & 0.00 [0.00, 0.00] & --- & --- \\
rag & 1/28 & 0.04 [0.00, 0.11] & +0.04 [+0.00, +0.11] & 1.00 \\
hybrid (kg+rag) & 12/28 & 0.43 [0.25, 0.61] & +0.43 [+0.25, +0.61] & 0.0005\textsuperscript{***} \\
\textbf{kg (deployed CPG)} & 15/28 & 0.54 [0.36, 0.71] & +0.54 [+0.36, +0.71] & $<$0.0001\textsuperscript{***} \\
kg + inheritance edges & 21/28 & 0.75 [0.57, 0.89] & --- & --- \\
kg idealised (ceiling) & 26/28 & 0.93 [0.82, 1.00] & +0.93 [+0.82, +1.00] & $<$0.0001\textsuperscript{***} \\
\bottomrule
\end{tabular}
\end{table*}


Fig.~\ref{fig:detection-ladder} and Table~\ref{tab:exp2-structural} carry the
paper's central result. The diff-only
baseline detects none of the 28 cross-file defects. This is not a poor score;
it is the expected score, and it is worth pausing on, because it means the
capability is absent rather than weak. Embedding retrieval detects one. Graph
context detects fifteen ($\Delta = +0.54$, $[+0.36, +0.71]$, $p<0.0001$),
exceeding the pre-registered threshold of $+0.30$.

The retrieval result deserves comment, because it is the sharpest contrast
with Experiment 1, where retrieval scored \emph{highest} on total rubric
score. Retrieval fails here for a principled reason rather than an
implementation one: a defect's dependent is not textually similar to the
defect. A function's callers do not resemble the function. Similarity search
therefore retrieves the wrong neighbourhood, competently. The hybrid
condition detects twelve, slightly \emph{below} the graph alone, consistent
with retrieval's chunks diluting attention on the graph block rather than
adding to it.

\subsection{The placebo band}
% Generated by scripts/generate_paper_tables.py — do not edit by hand.
% Source: experiments/2026-07-05_injection_exp2/out/results.json
\begin{table}[t]
\centering
\caption{Experiment 2, local-control bands ($n=12$): the placebo check. The
defect is visible inside the diff, so every arm should detect it and no arm
should gain from added context. All differences fall inside the
pre-registered parity window $[-0.15, +0.15]$.}
\label{tab:exp2-local}
\footnotesize
\setlength{\tabcolsep}{3.5pt}
\begin{tabular}{lccc}
\toprule
Arm & Det. & Rate & $\Delta$ vs base. [95\% CI] \\
\midrule
baseline & 11/12 & 0.92 & --- \\
rag & 12/12 & 1.00 & +0.08 [+0.00, +0.25] \\
hybrid & 11/12 & 0.92 & +0.00 [-0.25, +0.25] \\
kg & 10/12 & 0.83 & -0.08 [-0.25, +0.00] \\
kg + inherit. & 10/12 & 0.83 & --- \\
kg idealised & 10/12 & 0.83 & -0.08 [-0.25, +0.00] \\
\bottomrule
\end{tabular}
\end{table}


On the 12 local-control defects, every arm detects at a high rate --- the
defect is in the diff, as designed --- and the graph condition's difference
from baseline is $-0.08$ $[-0.25, +0.00]$, inside the pre-registered parity
window. Added context does not make the reviewer generically more vigilant.
Whatever produces $15/28$ on structural defects does not transfer to defects
where structure is irrelevant, which is the signature of a capability rather
than an artefact.

\subsection{Where the remaining misses come from}
The idealised graph detects $26/28$. Since it differs from the deployed
graph only in how the dependents were resolved, the gap between $15$ and
$26$ is attributable to graph construction rather than to the language
model's ability to use the graph. Concretely, the deployed builder either
missed relations or rendered them insalient.

One pre-declared band isolates a specific instance. In the six defects whose
consequence propagates through inheritance, adding import and inheritance
edges raises detection from $2/6$ to $5/6$ ($\Delta = +0.50$, meeting the
pre-registered $+0.30$ threshold). Across all 28 structural defects the
augmented builder reaches $21/28$ against $15/28$
($\Delta = +0.21$ $[+0.04, +0.39]$, $p = 0.071$).

The failures are also informative. The Java repository is the deployed
graph's weakest ($2/8$ structural), and in several of those misses the
correct edges \emph{were present} in the graph. Edge presence is therefore
necessary but not sufficient: how the neighbourhood is selected and rendered
matters as much as whether it was resolved.

\subsection{Judge validation}
Against the deterministic mention-oracle, judge agreement is 72\%
(gpt-4o-mini), 80\% (gpt-4o), and 93\% (gemini-2.5-flash). Critically, both
OpenAI judges produced \emph{zero} ``detected without mention'' flags: no
judge credited a detection the review did not actually contain. The panel's
errors are therefore conservative --- missed detections, not invented ones ---
so the reported detection rates are if anything lower bounds.

\section{Discussion}
\label{sec:discussion}

\subsection{Context type dominates context quantity}
The two experiments together support a claim neither supports alone. Graph
and retrieval context produce statistically indistinguishable aggregate
rubric scores and wildly different cross-file detection rates. If ``how much
context'' were the operative variable, the hybrid condition --- which has the
most --- would be best; it is not, on either the dependency subset in
Experiment 1 or structural detection in Experiment 2. What varies is whether
the retrieved context has the shape of the question being asked. Cross-file
impact is a reachability question, and reachability is what a graph encodes
and similarity does not.

\subsection{Implications for evaluating review assistants}
Our Experiment 1 is, by construction, a competent instance of the standard
protocol, and on its own it would have licensed a much weaker and partly
misleading conclusion: that context helps a little, and that the choice of
context type is a matter of taste. We would not have detected the $0/28$
versus $15/28$ gap. We take this as an argument for capability-targeted
evaluation with objective ground truth alongside rubric scoring, in the
spirit of Hong et al.'s critique of code review
benchmarks~\cite{hong2024deepcrceval}. The negative-control band and the
placebo band are both cheap to add and both did real work here.

\subsection{Implications for building them}
The ceiling result relocates the engineering problem. With a perfect
resolver the reviewing model finds $26/28$, so the bottleneck is not the
model's reasoning over graphs but the graph itself --- and, as the Java
failures show, the selection and rendering of the neighbourhood as much as
its construction. That is an encouraging division of labour, because static
resolution improves with ordinary engineering effort.

\section{Threats to Validity}
\label{sec:threats}

\textbf{Construct validity.} Experiment 1's outcome is an LLM panel's rubric
judgement, which inherits verbosity and specificity
biases~\cite{zheng2023judging} that its substantial inter-judge agreement
cannot rule out. We therefore treat its
aggregate effects as suggestive and rest the paper's claim on Experiment 2,
whose ground truth is static and whose panel we validated against a
mention-oracle. Experiment 2's outcome, detection of an injected defect, is
narrower than review quality; we claim the capability, not overall
usefulness.

\textbf{Internal validity.} Injected defects are synthetic and may be easier
or more stereotyped than real ones; we mitigated selection bias by choosing
sites direction-blind on structural fan-out before generating any review.
The builder-sensitivity comparison uses a different sample ($n=35$) and
rebuilt only the graph arm, so it supports only the baseline-versus-graph
contrast. Experiment 1's graph block is built by filename matching, so an
unknown fraction of the relations it renders do not exist; we did not measure
that fraction, and a precision estimate would strengthen the paper. A third
variant that resolves relations by parsing was also run, and reduced the
effect, but it regenerated its arms at a different generation temperature than
the headline, so it confounds builder with temperature and we do not report it
as a builder contrast. Generation temperature differs between the experiments ($0.0$ and
$0.3$) because they call different pipeline entry points; within each
experiment all arms share one setting, so no contrast is confounded, but the
two experiments' absolute scores are not directly comparable. No generation
seed is sent to the model on either path.

\textbf{External validity.} Five repositories and 40 pull requests in
Experiment 1, three repositories and 40 defects in Experiment 2, one
generator model (gpt-4o). The pre-registered target of 48 injections was not
met: two bands had insufficient supply (the operator band is Python-only,
and the Java repository has no barrel-file equivalent), and we report the
shortfall rather than back-filling it. One repository scope substitution was
documented before the run because of a frontend gap. Results may differ for
other languages, larger changes, and models with different context windows.

\textbf{Statistical validity.} With $n=40$ and $n=28$, confidence intervals
are wide, and several Experiment 1 effects sit near the significance
boundary; we report intervals throughout and avoid claiming precision we do
not have. Multiple conditions are compared against a shared baseline, so
individual $p$-values should be read as descriptive.

\section{Conclusion}
\label{sec:conclusion}

We compared graph-shaped and similarity-shaped repository context for LLM
code review under two designs. Measured by aggregate rubric score on 40 real
pull requests, the two are nearly interchangeable and both help modestly.
Measured by detection of defects whose consequences provably lie outside the
changed file, they are not interchangeable at all: $15/28$ against $1/28$,
against a diff-only baseline that detects none. A placebo band rules out
context volume as the explanation, and an idealised graph reaching $26/28$
places most of the remaining gap in graph construction rather than in the
model.

The practical recommendation follows from the gap between those two
measurements rather than from either one. If cross-file consequence is what
reviewers most need help with, then it should be measured directly, because
the aggregate score that a standard evaluation produces is nearly blind to
it.

% TODO before submission:
%  - Human study (6 pull requests, pairwise preference) is instrumented and
%    live but not yet analysed; add as a validation subsection or drop it
%    cleanly. Do not cite results that do not exist yet.
%  - Add missing bib entries flagged above; do not invent keys.
%  - Confirm venue, page budget, and author list.
%  - Consider a figure: the 0/28 vs 1/28 vs 15/28 vs 26/28 ladder is the
%    paper's headline and currently only appears as a table.

\section*{Data Availability}
All artefacts referenced here are produced by scripted pipelines: evidence
packs, the 160 generated reviews, per-judge votes, the injection manifest
with its resolved-dependent evidence lines, and the bootstrap and
permutation outputs. Every table in this paper is generated directly from
those artefacts by a single script, and an independent verification script
re-derives each reported number from the raw judge votes.
% TODO: replace with an anonymised or public repository link at submission.

\bibliographystyle{IEEEtran}
\bibliography{references}

\end{document}
